% !TEX program = lualatex
\documentclass[11pt, a4paper]{article}

\usepackage{fontspec}
\setmainfont{TeX Gyre Pagella}

\usepackage{amsmath, amssymb, amsthm}
\usepackage{geometry}
\usepackage{microtype}
\usepackage{setspace}
\usepackage{titlesec}
\usepackage{hyperref}
\usepackage{mdframed}

\geometry{
  top=1.25in,
  bottom=1.25in,
  left=1.35in,
  right=1.35in
}

\onehalfspacing

\hypersetup{
  colorlinks=true,
  linkcolor=black,
  citecolor=black,
  urlcolor=black
}

\newtheoremstyle{named}
  {12pt}{12pt}{\itshape}{}{\bfseries}{.}{.5em}{}
\theoremstyle{named}
\newtheorem{proposition}{Proposition}
\newtheorem{definition}{Definition}

\newmdenv[
  linewidth=0.8pt,
  innertopmargin=10pt,
  innerbottommargin=10pt,
  innerleftmargin=14pt,
  innerrightmargin=14pt,
  skipabove=14pt,
  skipbelow=14pt
]{keybox}

\titleformat{\section}{\large\bfseries}{\thesection.}{0.5em}{}
\titleformat{\subsection}{\normalsize\bfseries}{\thesubsection.}{0.5em}{}

\title{%
  \textbf{Too Coherent to Be Real}\\[0.4em]
  \large Szukalski, Creative Obstinacy, Documentary Unreality,\\
  and the Problem of the Frictionless Subject
}
\author{Flyxion}
\date{}

% ---------------------------------------------------------------
\begin{document}

\maketitle
\thispagestyle{empty}

\begin{abstract}
This essay begins with an unusual perceptual experience: watching a
documentary about the Polish sculptor Stanis\l{}aw Szukalski and being
left uncertain whether he actually existed. The film is so well-made,
its subject so perfectly illustrative of everything a documentary about
outsider genius, nationalist tragedy, and late redemption would want
him to be, that the viewer's reality-testing mechanism fires not on
any implausible detail but on the overall gestalt. Szukalski is too
coherent. His contradictions resolve too cleanly. The care with which
he has been preserved produces, paradoxically, a figure who feels
constructed rather than encountered. The essay uses this experience
as an entry point into three interconnected arguments. First, that
Szukalski's creative obstinacy — his systematic refusal of outside
influence — was not mere eccentricity but a survival strategy, a form
of procedural memory maintenance under conditions of institutional
destruction. Second, drawing on Catherine Liu's analysis of
infantilization, affect, and the professional managerial class, that
total mediation produces unreality: a subject with no friction, no
resistance to the narrative containing him, becomes indistinguishable
from a fabrication regardless of whether he was real. Third, that these
two observations are connected by a deeper structural claim: the
conditions that make creative work durable — obstinacy, refusal of
absorption, maintenance of idiosyncratic form — are also the conditions
that make a person difficult to document without flattening. The essay
concludes by reading Szukalski's Zermatism — his theory of a single
origin language and a debased lineage of humanity — as an armillary
proposal gone pathological: a coordinate transformation that kept
rotating past the point where it could encounter resistance from
reality, and thereby became unfalsifiable in the same way the egg-avatar
politics Liu diagnoses becomes unfalsifiable: sealed against the
friction that would have made it real.
\end{abstract}

\newpage
\tableofcontents
\newpage

% ---------------------------------------------------------------
\section{The Unreality Effect}

There is a specific perceptual experience that well-made documentaries
about eccentric geniuses can produce in the attentive viewer: the
growing suspicion that the subject did not exist. This is not the
ordinary skepticism one brings to historical claims. It is something
stranger. The figure on screen is too perfectly illustrative of
everything the film needs them to be. Their contradictions are too
narratively productive. Their tragedy is too clean. The people who
loved them remember them with exactly the combination of wonder,
sorrow, and complicated loyalty that the film requires. And gradually
the viewer's reality-testing mechanism, which usually fires on
implausible details, begins to fire instead on the overall coherence
of the construction. The subject is too much like what a subject of
this kind would be like.

Stanis\l{}aw Szukalski produces this effect. He is a Polish sculptor
born in 1893 who developed a distinctive and obsessive visual language,
was celebrated in interwar Poland, lost his entire Warsaw studio and
nearly his entire life's work to the Nazi bombing of 1939, emigrated
to Los Angeles, lived in obscurity for decades, and was rediscovered
in his old age by a circle of admirers including George DiCaprio, who
preserved his remaining work and recorded his conversation. He also
held nationalist and antisemitic views in the years surrounding the
Second World War, views he later repudiated. He created a theory of
human cultural origins he called Zermatism, which posited that all
human languages descended from a single primordial tongue he called
Protong, and that humanity's debasement arose from prehistoric
interbreeding with sea creatures.

Presented as a list of facts, this is simply unusual. Presented as a
documentary, with archival photographs, interviews with devoted
survivors, and footage of the actual sculptures — gnarly, mythological,
technically extraordinary — the unreality effect intensifies rather
than diminishes. The more evidence accumulates, the more the figure
feels assembled. He is too complete. The tragedy of the lost Warsaw
work, the obscurity, the late rescue, the moral complexity of the
nationalism — each element fulfills a narrative function so precisely
that the whole begins to feel authored.

This essay begins with that perceptual experience because it is
philosophically interesting in itself, and because it leads directly
into questions that are larger than any one documentary: questions
about creative obstinacy as a survival strategy, about the relationship
between documentation and reality, about what kinds of subjects resist
the flattening that total mediation produces, and about when the
refusal of outside influence becomes not a condition of creative
survival but a sealed system that can no longer encounter the world
on its own terms.

% ---------------------------------------------------------------
\section{Szukalski's Obstinacy as Survival Strategy}

The standard reading of an outsider artist's refusal of influence is
psychological: eccentricity, grandiosity, incapacity for the ordinary
social negotiations that participation in an artistic community
requires. This reading is not entirely wrong in Szukalski's case.
He was by all accounts difficult, domineering, convinced of his own
exceptional status in ways that made sustained collaboration nearly
impossible.

But the psychological reading misses the structural function of the
refusal. Szukalski's obstinacy — his insistence on developing a
visual and theoretical vocabulary entirely his own, his resistance to
assimilation into any existing artistic movement, his cultivation of
a mythology of himself as the lone bearer of a suppressed primordial
tradition — was not only a personality trait. It was a form of
institutional independence. And institutional independence, as the
history of the twentieth century repeatedly demonstrated, was a
condition of survival.

An artist whose work is legible within an existing movement,
marketable through existing institutional channels, comprehensible to
existing critical frameworks, is also an artist whose work is
vulnerable to the institutional forces that shape and destroy those
frameworks. The Warsaw bombing of 1939 destroyed Szukalski's studio,
his archive, and the physical record of decades of work. What it could
not destroy was the procedural knowledge carried in his hands and his
mind: the specific techniques, the mythological vocabulary, the
obsessive formal investigations that constituted his practice. That
knowledge survived because it was his alone. There was no school of
Szukalski, no workshop tradition, no institutional structure through
which the work was mediated and which could have been infiltrated or
corrupted or simply closed.

\begin{keybox}
The artist who refuses absorption preserves something that institutions
cannot destroy, because institutions cannot absorb what they cannot
recognize. Szukalski's idiosyncrasy was not an obstacle to his
survival. It was its mechanism.
\end{keybox}

This is the inverse of the standard account of institutional memory.
In most analyses, procedural knowledge is preserved through
transmission — through workshops, apprenticeships, training
programs, communities of practice. When those transmission mechanisms
are destroyed, the knowledge is lost. Szukalski's case suggests a
complementary but different dynamic: knowledge can also be preserved
through radical concentration in a single carrier who is too
idiosyncratic to be absorbed by any institution and therefore cannot
be destroyed when institutions are destroyed.

The cost of this strategy is obvious. Knowledge preserved through
radical concentration in a single carrier is also knowledge that
cannot be transmitted through normal channels, that depends entirely
on the carrier's survival and willingness to continue, and that is
vulnerable to exactly the death and dispersal that eventually claimed
Szukalski's unreproduced work. The DiCaprio circle's preservation
effort is, in the framework developed in recent work on institutional
memory, an attempt at approximate repair under conditions of
operational fiber death: the procedural knowledge that was Szukalski's
alone has decayed to measure zero, and what remains is declarative
knowledge — photographs, recordings, the surviving sculptures — from
which the full practice cannot be reconstructed, only imagined.

% ---------------------------------------------------------------
\section{The Egg Avatar and Its Inverse}

Catherine Liu, in conversation about her work on the professional
managerial class and the culture of liberal affect, describes a figure
she calls the egg avatar: the undifferentiated online subject who has
cultivated maximum receptivity to feeling at the cost of the ego
boundaries that would allow actual confrontation with the world. The
egg avatar is soft, porous, constantly vibing, constitutionally
incapable of the friction that distinguishes real encounter from
mutual affect-resonance. When egg avatars collide, they do not
disagree; they fear cracking. The political consequence is an inability
to build anything that requires sustained conflict, institutional
structure, or the tolerance of contradiction.

Liu is drawing on a psychoanalytic tradition that distinguishes between
the ego — the mediating structure between id and superego, capable of
engaging with external reality — and the pre-differentiated infantile
surface that precedes ego formation. Her argument is that a significant
strand of contemporary progressive culture has sacralized the
pre-differentiated state, treating the porousness and receptivity of
the undifferentiated infant as a political and aesthetic ideal rather
than a developmental starting point to be moved through.

Szukalski is the opposite of an egg. He is what happens when
differentiation proceeds without external check — when the ego boundary
is maintained and reinforced and elaborated until it becomes a
totalizing system that can no longer register feedback from outside
itself. If the egg avatar's pathology is the inability to sustain
boundaries, Szukalski's pathology is the inability to relax them.
Where the egg avatar dissolves into collective affect, Szukalski
constructed an entire cosmology as an extension of his own
self-understanding. Zermatism is not a theory that could be falsified
by evidence because it was not primarily a theory about the world; it
was a theory about Szukalski's position in relation to the world, and
evidence that failed to confirm it was evidence of the debased
condition the theory was designed to diagnose.

\begin{keybox}
The egg avatar and the totalizing obsessive are not opposites in the
sense of correcting each other. They are mirror pathologies of the
same failure of mediation: one has no boundary, the other has a
boundary that has become a wall.
\end{keybox}

Liu's framework implies that what is needed is not more differentiation
or less differentiation but the right kind: the ego as mediating term,
capable of sustaining conflict without dissolution, capable of
encountering the world without becoming the world. Szukalski had
enormous ego strength in the psychoanalytic sense — he sustained his
practice through decades of obscurity, institutional destruction, and
exile — but his ego had long since ceased to mediate between his inner
world and external reality. It had become, instead, a system for
interpreting external reality as confirmation of an already-complete
internal theory.

This is where his nationalism and antisemitism cease to be incidental
biographical facts and become structurally interesting. A mind that
has developed a totalizing cosmological system — Zermatism, Protong,
the theory of prehistoric debasement — is a mind that has solved the
problem of outside influence by eliminating the category of outside.
Everything becomes interpretable within the system. The antisemitism
is not a separate module added to an otherwise coherent worldview; it
is what a Zermatist interpretation of early twentieth century European
politics produces when the interpretive frame has become sealed.
Szukalski did eventually repudiate these views, which suggests the
frame was not entirely impermeable. But the repudiation came late
and appears to have been experienced as a correction from outside
rather than an internal revision — the kind of encounter with
recalcitrant reality that a more porous ego structure would have
enabled earlier.

% ---------------------------------------------------------------
\section{What the Documentary Cannot Show}

The documentary about Szukalski is made by people who love him.
Glenn Bray, who preserved the work; George DiCaprio, who befriended
him in his Los Angeles years; others who experienced his personality
directly and who carry genuine grief about what was lost in Warsaw.
Their testimony is authentic. Their care for the man and his work is
unmistakable.

And yet this care is precisely what produces the unreality effect.
Total mediation through love creates a subject without friction.
The Szukalski we see in the documentary has been selected, arranged,
and preserved by people whose relationship to him was one of devotion.
His difficult qualities — the domineering personality, the grandiosity,
the nationalism — are present but contextualized, explained,
integrated into a narrative of complexity-that-is-ultimately-redeemed.
The result is not dishonest. It is something more interesting: it is
a documentary that has solved all the interpretive problems its subject
poses, and in solving them has made the subject less real.

Liu's analysis of affect is useful here. The film produces affect in
the viewer: wonder at the sculptures, discomfort at the nationalism,
sympathy for the losses, satisfaction at the late recognition. These
are genuine responses to genuine material. But affect, in Liu's
framework, is not the same as intersubjectivity. Affect flows; it
connects surfaces. Intersubjectivity requires differentiation, the
acknowledgment that the other is genuinely other — not reducible to
what they mean to me, not fully legible within my interpretive
framework, capable of being wrong in ways I did not anticipate.

The documentary gives us a Szukalski who is fully legible. His
contradictions are there but they have been made narrative. He is
not, in the viewing experience, genuinely other. He is an
extraordinarily well-realized figure in someone else's story about
him. And this is the source of the unreality effect: a person
encountered through total mediation by devoted interpreters is a
person whose alterity has been managed. He exists inside the love
of the people who preserved him, and inside that love he is as
complete and as real as a character in a very good novel.

The doubt is not about whether Szukalski existed. It is about
whether the documentary's Szukalski and the historical Szukalski
are the same entity. And that doubt is structurally correct:
they are not. The historical Szukalski produced friction — with
institutions, with collaborators, with political movements, with
the limits of his own cosmological system. The documentary's Szukalski
produces affect. These are different modes of being in the world,
and the documentary, by excellence of its craft, has substituted one
for the other.

% ---------------------------------------------------------------
\section{The Observer Problem and the Chain of Love}

The UPC framework developed by Eloy Escagedo Gutierrez argues that
paradoxes arise when observers are removed from explanatory
frameworks. There is a related but inverse problem: what happens when
observers are so present, so constitutive of the object being
described, that the object cannot be encountered independently of
them?

The Szukalski documentary is organized around a chain of observers
who are also participants: the people who knew him, loved him, grieved
for him, and devoted significant portions of their lives to preserving
what he made. This chain of love is also a chain of mediation. Each
link in the chain — the photographs Bray collected, the conversations
DiCaprio recorded, the documentary Irek Dobrowolski assembled from
these materials — is an act of preservation that is also an act of
transformation. What is preserved is not Szukalski but what Szukalski
meant to the people who preserved him.

This is not a criticism of the film. It is a description of what
documentation always does, made visible by the particular intensity
of devotion in this case. All archives are partial. All preservation
is selection. All documentation reflects the purposes and frameworks
of the documenters. But usually these purposes and frameworks are
diverse enough that the subject partially escapes them — some witnesses
loved the subject, some were indifferent, some were hostile, and the
friction between their accounts produces something like an independent
object that the various accounts are inadequately describing.

When all the witnesses loved the subject, that friction disappears.
The subject becomes the consensus object of devotion, and the
consensus object of devotion is, by definition, what the devotees
needed it to be. Szukalski's reality — his specific recalcitrance,
the ways in which he was not what anyone needed him to be — is
present in the historical record but has been largely smoothed into
the documentary's arc of tragic greatness and belated recognition.

\begin{keybox}
The observer cannot be removed from the chain of meaning. But when
the chain of meaning is also a chain of love, the observer's
presence becomes so total that the observed disappears into the
observation.
\end{keybox}

What the documentary cannot show is the Szukalski who was difficult
in ways that were not productive, wrong in ways that were not
redeemable, ordinary in ways that were not interesting. It cannot
show the texture of his actual daily obstinacy — the meetings that
went nowhere, the collaborations that failed for mundane rather than
tragic reasons, the conversations that were simply boring or
repetitive or self-absorbed. These are the friction-producing elements
of a real person, and they are exactly what love, in the act of
preservation, edits out.

% ---------------------------------------------------------------
\section{Zermatism as an Armillary Proposal Gone Sealed}

Szukalski's theoretical system — Zermatism, with its claims about
Protong as the single origin language, its theory of prehistoric
interbreeding, its morphological analyses of cultural symbols across
unrelated civilizations — can be read as an armillary proposal: a
coordinate transformation applied to the questions of human cultural
origin, linguistic diversity, and civilizational development.

Read charitably and structurally, Zermatism asks a legitimate
question. What if the diversity of human languages and cultural forms
is not a branching process of divergence from a common point but a
pattern of corruption and fragmentation from a unified original? What
if the resemblances between morphologically distant linguistic and
cultural systems are not coincidental but trace to a common source
that standard comparative linguistics has failed to recover? What if
the history of civilization is better understood as a story of loss
than of development?

These are not crazy questions. Versions of them appear in serious
comparative mythology, in the study of linguistic universals, in
debates about the deep history of human symbolic culture. Zermatism
is what happens when these questions are pursued by a mind that has
lost the mechanism for encountering disconfirming evidence — when the
armillary proposal keeps rotating past the point where it can be
checked against the territory it is supposed to be mapping.

The armillary sphere is a navigation instrument. Its value depends on
its relationship to the celestial mechanics it represents. An armillary
sphere that has been decoupled from astronomical observation — that
has become self-referential, its rings encoding relationships that
are internally consistent but no longer checked against the sky —
ceases to be a navigation instrument and becomes a beautiful closed
object. It still looks like a sphere. It still has rings. But it has
lost the property that made it useful: its capacity to be wrong.

Zermatism had lost that capacity. It could not be wrong because
every piece of evidence that might have falsified it was
reinterpreted as evidence of the debasement and confusion it was
designed to diagnose. Morphological similarities between distant
languages confirmed Protong. Morphological differences were evidence
of corruption. The theory was sealed.

This is the pathological form of the armillary proposal, and it is
structurally identical to what Liu diagnoses in a different register
when she describes the egg avatar's relationship to political
discourse. The egg avatar cannot be wrong either — not because its
theory is too powerful but because its porousness prevents the
formation of the boundaries that would allow genuine disagreement.
One has no boundary; the other has a boundary that has become
impermeable. Both are sealed against the encounter with recalcitrant
reality that is the condition of genuine knowledge.

The connection to the antisemitism is direct and uncomfortable.
A cosmological system that posits a debased lineage contaminating
pure human stock, developed in interwar Europe by a Polish nationalist,
does not require external ideological influence to arrive at
antisemitic conclusions. It arrives there through the internal logic
of the system once the system has been decoupled from the capacity
to encounter disconfirming evidence. The antisemitism is not a foreign
body in Zermatism. It is what Zermatism looks like when sealed.

% ---------------------------------------------------------------
\section{The PMC and the Institutionalization of the Sealed System}

Liu's analysis of the professional managerial class is relevant to
Szukalski in an unexpected way. Her central argument is that the PMC
has substituted the performance of virtue for the exercise of reason —
that it has developed a doctrinal system (identity politics,
decolonization frameworks, language rituals) that functions as a
closed interpretive system: internally consistent, resistant to
external critique, organized around excommunication of doctrinal
error rather than engagement with contradicting evidence.

The PMC's sealed system is the institutional form of what Szukalski's
sealed system was at the individual level. Both are products of
differentiation that has proceeded past the point of productive
encounter with the world. The PMC variant is more dangerous in some
ways because it has institutional infrastructure — academic positions,
editorial boards, grant committees, human resources departments —
through which the sealed system can be enforced on people who did not
choose to inhabit it.

But the structural similarity runs deeper than this. Liu argues that
the PMC's virtue culture is crypto-religious: it has the form of
doctrinal systems, with heresy, excommunication, and the performance
of correct belief as markers of membership, without the substantive
theological content that would anchor the doctrine to anything
outside itself. Zermatism has the same structure at the individual
level: it is a total explanatory system that accounts for all evidence
and is therefore unfalsifiable, organized around a mythology of
primordial purity and subsequent corruption that gives its adherent
a position of interpretive superiority over those who do not share
the framework.

Both systems produce a similar pathology in their adherents: the
inability to learn from encounter with the world, because the world
has been pre-interpreted by the system. Both systems also produce,
paradoxically, a kind of passionate commitment that appears to
outsiders as evidence of genuine engagement. The PMC activist who
has spent years developing fluency in the doctrinal vocabulary of
their movement looks, from inside the movement, like someone who
cares deeply. From outside, they look like someone who has
substituted the performance of caring for the experience of caring
— exactly the distinction Liu draws between genuine maternal
attentiveness and its liberal dilution into performative concern.

\begin{keybox}
Sealed systems produce intense internal commitment and impermeability
to external reality simultaneously. The intensity is real. It is also
exactly what prevents the system from doing the work it claims to do.
\end{keybox}

What both Szukalski and the PMC demonstrate, from different angles,
is that differentiation without the capacity for genuine encounter
with recalcitrant reality produces closure rather than integrity.
Ego strength that cannot be surprised is not strength; it is rigidity.
A cosmological system that cannot be falsified is not a theory; it
is a monument to its own construction.

% ---------------------------------------------------------------
\section{Obstinacy, Maturity, and the Problem of the Canon}

Liu argues, following Kant, that democratic citizenship requires
maturity — the willingness to use one's own reason, to emerge from
self-imposed tutelage, to engage with the world as it is rather than
as one's affect-resonance with a group suggests it should be. She
notes that contemporary academic culture has largely inverted this
value, treating maturity as normative, discipline as oppressive, and
the cultivation of permanent infantile receptivity as a form of
liberation.

Szukalski's case complicates this picture in a productive way. He
achieved something like the opposite of infantile receptivity: a
radically individuated creative practice that developed entirely on
its own terms, refused absorption into existing movements, and
maintained its procedural integrity through decades of institutional
destruction and personal obscurity. By the standards Liu is applying,
he had strong ego, high differentiation, and a sophisticated capacity
for sustained intellectual and creative labor. He was, in the
psychoanalytic sense, mature.

And yet the maturity had curdled. The differentiation had become
enclosure. The refusal of outside influence, which was the condition
of his survival, had also become the condition of his antisemitism
and his unfalsifiable cosmology. The strong ego, unmediated by genuine
encounter with resistant otherness, had produced a totalizing system
rather than a well-developed self.

This suggests that Liu's account, while correct about the pathology
of the egg avatar, needs a complementary account of the pathology
of excessive differentiation. Maturity is not simply the achievement
of ego strength and differentiation. It requires, in addition, the
maintenance of genuine openness to encounter — the willingness to
be surprised by the world, to be wrong, to revise one's frameworks
in response to evidence that does not fit them.

The canon that Liu defends — the tradition of Enlightenment reason,
scientific method, Marxist political economy, psychoanalytic
self-understanding — is valuable not merely because it represents
accumulated knowledge but because it represents accumulated failure.
The canon has been wrong. It has been revised. It contains arguments
that were superseded, predictions that did not pan out, frameworks
that proved inadequate to the phenomena they were designed to explain.
This revisability is precisely its value: it is a record of genuine
encounter with recalcitrant reality. The Enlightenment tradition did
not produce a sealed system. It produced a practice of self-correction.

Szukalski never developed a practice of self-correction. His late
repudiation of his nationalist views suggests he was capable of it,
but by then the cosmological system that had licensed those views
had been in place for decades, and the repudiation appears to have
been a single correction rather than the beginning of a revisionary
practice. Zermatism survived his moral revision intact.

% ---------------------------------------------------------------
\section{Operational Fiber Death and the Lost Warsaw Work}

The destruction of Szukalski's Warsaw studio in 1939 is an instance
of what the institutional attractors framework calls operational fiber
death: the condition in which declarative knowledge of a prior
configuration survives while the procedural capacity to reconstruct
it decays to measure zero.

The photographs of the Warsaw work exist. Descriptions of it exist.
The surviving sculptures provide evidence of the formal vocabulary
and the technical methods. But the Warsaw works themselves are gone,
and the specific procedural knowledge that produced them — the
particular solutions to particular formal problems, the decisions made
in the making that are not visible in the finished object — decayed
with Szukalski's death. What the DiCaprio circle preserved was the
declarative memory. The procedural memory was not transmissible
because Szukalski never transmitted it, and when he died it died
with him.

This is the irony at the center of his survival strategy. The
obstinacy that preserved him through institutional destruction was
also the obstinacy that prevented the transmission of his practice
to anyone who could have carried it forward. His work survived the
bombing because it was entirely his; it could not survive his death
for the same reason.

The documentary is itself an exercise in approximate repair: an
attempt to recover function — the experience of Szukalski's work,
the sense of his creative intelligence, the testimony of people who
encountered it — without recovering form. The surviving sculptures
are not the lost Warsaw sculptures. The recordings of his
conversation are not his creative practice. The documentary is a
functional equivalent of having known him, not a restoration of what
he actually was.

Approximate repair is what is available when operational fiber death
has occurred. The framework developed elsewhere argues that this is
a correction of the target for reformers: the appropriate goal is
not restoration but functional reconstruction, the design of new
forms capable of recovering lost functions in the current environment.
Applied to Szukalski, this means that the appropriate relationship
to his work is not curatorial — the attempt to preserve and transmit
what he made as accurately as possible — but generative: the attempt
to understand what questions his work was probing and to continue
probing those questions through new work.

Whether the documentary achieves this is uncertain. It achieves
something more modest: it makes visible that something was lost,
and makes that loss felt. That is not nothing. Making a loss visible
is the first step toward understanding what kind of repair might be
possible.

% ---------------------------------------------------------------
\section{Reachability and the Sealed Cosmology}

The reachability framework developed in the institutional attractors
analysis asks: what futures remain accessible from the current
configuration? The parallel question for a creative and intellectual
practice is: what revisions remain accessible from within the current
framework?

For Szukalski, the answer by the middle of his career was: very few.
Zermatism had become a sealed interpretive system that could
accommodate any new evidence by reinterpreting it as confirmation of
the theory. A practice organized around a sealed system cannot be
revised from within; it can only be abandoned or maintained. Szukalski
maintained it.

The cost of the sealed system is not merely that it produces wrong
conclusions. It is that it progressively eliminates the futures that
would have been accessible from a more porous framework. Every year
that Zermatism was maintained without revision was a year in which
the possibility of a different Szukalski — one who had engaged with
comparative linguistics, with archaeological evidence, with the
intellectual traditions that might have challenged and enriched his
morphological speculations — became less reachable. The admissible
set of creative futures contracted as the system sealed.

This is the precise sense in which Szukalski's obstinacy, which was
a condition of his survival, was also a condition of his limitation.
The institutional independence that preserved him from external
destruction also prevented the internal revision that would have kept
his intellectual practice alive as a practice rather than a monument.
He survived. His practice did not survive him. And the practice did
not survive him partly because it had, long before his death, ceased
to be a practice in the relevant sense: a mode of inquiry capable of
producing surprise.

% ---------------------------------------------------------------
\section{The Tyranny of Affect and the Tyranny of System}

Liu describes the tyranny of affect as the substitution of feeling
for objectivity in political and cultural discourse: the claim that
shared affect constitutes shared reality, that the intensity of
collective feeling licenses conclusions that evidence and argument
could not support, that skepticism and dissent are forms of violence
against the affective community rather than contributions to a shared
project of understanding.

Szukalski enacted a different but structurally related tyranny: the
tyranny of system. Where the tyranny of affect refuses objectivity
by dissolving the boundary between the feeling subject and the world,
the tyranny of system refuses objectivity by totalizing it — by
constructing a system so comprehensive that every piece of the world
is already interpreted before it is encountered. Both are forms of
impermeability to reality. Both produce, in their adherents, a
specific combination of passionate commitment and invulnerability to
evidence.

The political consequence of the tyranny of affect, in Liu's analysis,
is the inability to build organizations capable of sustained conflict
with actual power. Eggs cannot win. They crack against each other
before they reach the forces that are actually organizing the world.
The political consequence of the tyranny of system is different but
equally disabling: it produces certainty without knowledge, action
without feedback, and eventually the kind of moral catastrophe that
a fully sealed interpretive system licenses when it encounters actual
political conditions, as Zermatism encountered the actual political
conditions of interwar Poland.

What Liu's account implies, and what Szukalski's career illustrates
from the other direction, is that the problem of political and
creative agency is a problem of the right kind of boundary. Too soft:
dissolution into collective affect, inability to sustain conflict,
the egg-avatar politics that cannot build anything. Too hard:
enclosure into a totalizing system, inability to learn from
recalcitrant evidence, the sealed-cosmology politics that can build
things but cannot correct them.

The mediating term — the ego in the psychoanalytic vocabulary, reason
in the Enlightenment vocabulary, the practice of self-correction in
the scientific vocabulary — is what allows a boundary to remain a
boundary without becoming a wall. Szukalski had the boundary without
the self-correction. The PMC has institutional structures without the
ego functions those structures require to operate as something other
than enforcement mechanisms. The egg avatar has neither boundary nor
structure.

\begin{keybox}
What is needed is not more differentiation or less differentiation
but the capacity to hold a boundary that can also be revised — an
ego that is strong enough to sustain encounter with the world and
porous enough to be changed by it.
\end{keybox}

% ---------------------------------------------------------------
\section{The Documentary as Functional Equivalent}

What, finally, should we make of the documentary that produced the
unreality effect? The film is, by any ordinary measure, excellent.
The sculptures are extraordinary. The testimony of the people who
knew Szukalski is moving. The account of the Warsaw loss is genuinely
terrible. The navigation of his nationalism is careful and honest
without being either exculpatory or simply condemnatory.

And yet the unreality effect is a legitimate response to the film,
and understanding why it arises is more useful than dismissing it
as uncharitable to the filmmakers. The unreality effect is a signal
that something has been managed, that the friction-producing elements
of a genuinely recalcitrant subject have been processed into
narrative and affect. Szukalski as encountered in the documentary is
a functional equivalent of Szukalski — it recovers the function (the
experience of his intelligence, the sense of his loss, the encounter
with his work) without recovering the form (the actual texture of
his difficult, wrong, obstinate, sealed, sometimes terrible and
sometimes extraordinary mind).

This is what approximate repair looks like in documentary form. It
is not dishonest. It is the best available approximation given that
operational fiber death has occurred — given that the procedural
knowledge of what it was actually like to be in a room with him, to
encounter his thinking before it had been curated by love, to
experience the full weight of his sealed system's implications, is
not recoverable. The documentary gives us what it has, and what it
has is genuine and valuable.

The unreality effect is what real people feel like when they have been
loved into narrative coherence. It is not a failure of the documentary.
It is evidence that Szukalski was real — that he was recalcitrant
enough to produce, even through the total mediation of devoted
preservation, the sense that something has escaped the frame.

Something has. That is what reality does.

% ---------------------------------------------------------------
\section*{Conclusion}

Stanis\l{}aw Szukalski was too coherent to be real and real enough to
exceed his documentation. His creative obstinacy was a survival
strategy that worked and a closure mechanism that sealed him. His
cosmological system was an armillary proposal that kept rotating past
the point of encounter with the territory it claimed to map. His
late rediscovery was an act of love that produced, through the
excellence of its preservation, a functional equivalent of a man
who could not be fully recovered.

Liu's framework gives us the conceptual vocabulary for understanding
why: the egg avatar and the sealed system are mirror pathologies of
the same failure of mediation, and what is needed in both cases is
the ego function — the boundary that can be revised, the
differentiation that remains permeable to evidence, the maturity
that knows itself without sealing itself against the world.

The documentary's unreality effect is the viewer's correct intuition
that they are encountering a functional equivalent rather than the
thing itself. The thing itself is gone. What remains is the
documentation of the love that preserved what could be preserved,
and the sculptures, which still produce the friction that love alone
cannot manufacture.

The sculptures are real. They push back. They do not resolve into
narrative. They are too strange to be fully explained by any account
of the man who made them, including his own.

That is the only kind of reality that survives total mediation:
the kind that continues to resist being fully known.

% ---------------------------------------------------------------
\begin{thebibliography}{99}

% On Szukalski
\bibitem{szukalski1923}
Szukalski, S. (1923).
\textit{Works of Szukalski}.
University of Chicago Press.

\bibitem{bray2000}
Bray, G. (Ed.). (2000).
\textit{Behold!!! The Protong}.
Last Gasp.

% Documentary and mediation
\bibitem{nichols2001}
Nichols, B. (2001).
\textit{Introduction to Documentary}.
Indiana University Press.

\bibitem{renov1993}
Renov, M. (Ed.). (1993).
\textit{Theorizing Documentary}.
Routledge.

\bibitem{corner1996}
Corner, J. (1996).
\textit{The Art of Record: A Critical Introduction to Documentary}.
Manchester University Press.

% Psychoanalysis and ego theory
\bibitem{freud1923}
Freud, S. (1923).
\textit{Das Ich und das Es}.
Internationaler Psychoanalytischer Verlag.
[English: \textit{The Ego and the Id}. Norton, 1962.]

\bibitem{hartmann1958}
Hartmann, H. (1958).
\textit{Ego Psychology and the Problem of Adaptation}.
International Universities Press.

% Liu and PMC analysis
\bibitem{liu2021}
Liu, C. (2021).
\textit{Virtue Hoarders: The Case Against the Professional Managerial Class}.
University of Minnesota Press.

\bibitem{ehrenreich1977}
Ehrenreich, B., \& Ehrenreich, J. (1977).
The Professional-Managerial Class.
\textit{Radical America}, 11(2), 7--31.

% Enlightenment, maturity, reason
\bibitem{kant1784}
Kant, I. (1784).
Beantwortung der Frage: Was ist Aufkl\"{a}rung?
\textit{Berlinische Monatsschrift}.
[English: What is Enlightenment?]

\bibitem{habermas1981}
Habermas, J. (1981).
\textit{Theorie des kommunikativen Handelns}.
Suhrkamp.
[English: \textit{The Theory of Communicative Action}. Beacon Press, 1984.]

% Affect theory and its critics
\bibitem{massumi2002}
Massumi, B. (2002).
\textit{Parables for the Virtual: Movement, Affect, Sensation}.
Duke University Press.

\bibitem{leys2011}
Leys, R. (2011).
The Turn to Affect: A Critique.
\textit{Critical Inquiry}, 37(3), 434--472.

% Organizational structure and the left
\bibitem{freeman1972}
Freeman, J. (1972).
The Tyranny of Structurelessness.
\textit{Berkeley Journal of Sociology}, 17, 151--164.

% Outsider art and institutional independence
\bibitem{cardinal1972}
Cardinal, R. (1972).
\textit{Outsider Art}.
Studio Vista.

\bibitem{maizels1996}
Maizels, J. (Ed.). (1996).
\textit{Raw Creation: Outsider Art and Beyond}.
Phaidon.

% Memory, loss, and institutional decay
\bibitem{connerton1989}
Connerton, P. (1989).
\textit{How Societies Remember}.
Cambridge University Press.

\bibitem{nora1989}
Nora, P. (1989).
Between Memory and History: Les Lieux de M\'{e}moire.
\textit{Representations}, 26, 7--24.

% Falsifiability and closed systems
\bibitem{popper1959}
Popper, K. R. (1959).
\textit{The Logic of Scientific Discovery}.
Hutchinson.

\bibitem{kuhn1962}
Kuhn, T. S. (1962).
\textit{The Structure of Scientific Revolutions}.
University of Chicago Press.

% Nationalism and creative thought
\bibitem{anderson1983}
Anderson, B. (1983).
\textit{Imagined Communities: Reflections on the Origin and Spread of Nationalism}.
Verso.

% Comparative mythology and linguistics
\bibitem{dumezil1958}
Dum\'{e}zil, G. (1958).
\textit{L'Id\'{e}ologie tripartite des Indo-Europ\'{e}ens}.
Latomus.

\bibitem{campbell1949}
Campbell, J. (1949).
\textit{The Hero with a Thousand Faces}.
Pantheon.

\end{thebibliography}

% ---------------------------------------------------------------
\end{document}
